\subsection{Group Size, Consensus Quality, and Between-Team Variation}
\label{sec:group-size}

Adding annotators transforms the quality distribution in two distinct stages. On the fixed panel of 171 Manual panoramas, increasing team size from 3 to 6 raises mean $Q$ by 4.739 points (95\% path-bootstrap CI $[4.629,4.851]$). From 20 to 30, the mean gains only 0.500 points, whereas the within-image 10th percentile gains 2.848 (Figure~\ref{fig:group-size-171}, Table~\ref{tab:group-size-intervals}). Later additions primarily lift lower-tail quality and narrow differences between sampled teams. Evaluating group size solely by mean quality therefore misses this late-stage stabilization.

\begin{figure*}[t]
\centering
\setcounter{subfigure}{0}
\begin{minipage}{\textwidth}
\centering
{\footnotesize
\textcolor{figAll}{\rule[2pt]{10pt}{1.2pt}} All (171)\quad
\textcolor{paperBlue}{\rule[2pt]{10pt}{0.7pt}} Easy (116)\quad
\textcolor{figGreen}{\rule[2pt]{4pt}{0.7pt}\hspace{2pt}\rule[2pt]{4pt}{0.7pt}} Moderate (28)\quad
\textcolor{paperOrange}{\rule[2pt]{6pt}{0.7pt}\hspace{2pt}\rule[2pt]{1pt}{0.7pt}} Difficult (27)\par}
\vspace{3pt}
\ESPanel{.2414\linewidth}{panels/fig10/a.pdf}{Mean Q}%
\hfill%
\ESPanel{.2414\linewidth}{panels/fig10/b.pdf}{Lower-tail Q (P10)}%
\hfill%
\ESPanel{.2414\linewidth}{panels/fig10/c.pdf}{Between-team Q variation}%
\hfill%
\ESPanel{.2414\linewidth}{panels/fig10/d.pdf}{Complete-and-evaluable rate}

\end{minipage}
\caption{Group size affects mean Q, lower-tail Q, between-team variation and output completion differently. Curves use the same 171 images: 116 easy, 28 moderate and 27 difficult. Q summaries use scorable outputs; complete-and-evaluable rates use all 43,776 attempts per size. P10 is the image-specific 10th percentile, averaged over images. Shading marks 3--6 and 20--30 people. The at-least-half threshold admits exactly half of an even-sized team: the required support stays at 2 votes from 3 to 4 people, then rises to 3 at 5 people, contributing to odd--even alternation.}
\label{fig:group-size-171}
\end{figure*}


\begin{table}[t]
\centering\small
\caption{Changes from 20 to 30 people on the fixed 171-image panel, with equal image weights. Supplement D2 reports the quality levels.}
\label{tab:group-size-milestones}
\label{tab:group-size-intervals}
\setlength{\tabcolsep}{3pt}
\begin{tabular}{@{}lrrr@{}}
\toprule
Outcome & Change & Path interval & Building interval\\
\midrule
Mean Q & +0.500 & [0.439, 0.563] & [0.037, 1.042]\\
P10 & +2.848 & [2.620, 3.136] & [1.942, 3.970]\\
SD of Q & $-$2.031 & [$-$2.087, $-$1.962] & [$-$2.638, $-$1.503]\\
\bottomrule
\end{tabular}
\par\smallskip
\parbox{\columnwidth}{\footnotesize Intervals are approximate 95\% intervals. Path resampling measures Monte Carlo precision within the observed pools; building resampling retains the path estimates and resamples 22 buildings.}
\end{table}

The small late-stage mean gain also masks substantial image-level shifts. Among 27 difficult panoramas, 16 improve and 11 deteriorate from 20 to 30 people. Their signed mean change is $-0.174$ points, whereas their mean absolute change reaches 4.772. Opposing shifts cancel in the aggregate. Added observations alter both corner support and median positions: a turn with 4 supporters meets the retention threshold in an 8-person team but falls below it when a ninth member omits it. New endpoints also enter the median calculation, changing the location of retained turns when the median shifts. A flat average quality curve therefore conceals changes in disputed details.

Completion and quality measure distinct aspects of aggregation. The complete-and-evaluable rate fluctuates at small team sizes: it drops from 98.652\% at 3 people to 96.795\% at 6 (a net decrease of 1.857 percentage points) before reaching 99.488\% at 30. Quality evaluates scorable layouts; completion tracks whether an attempted team yields a complete, scorable layout. The 43,212 paths scorable at both 20 and 30 people preserve the same directions of mean, lower-tail and dispersion changes (Supplement D2). Separating these two metrics prevents conflating geometric accuracy with output validity.

Finite candidate pools contribute to this narrowing variation because larger teams share more members. Excluding the 12 images whose 30-person teams exhaust the pool leaves a P10 gain of 2.942 points across 159 images, demonstrating that lower-tail gains persist alongside growing member overlap (Supplement C10). In practice, late-stage additions mainly enforce consistency across teams drawn from finite annotator pools.
